\chapter{Att kontrollera ett påstående}
\label{app:verifiera}

Det här materialet gör, som allt undervisningsmaterial, en rad
påståenden.  De flesta går att pröva vid datorn, och det är hela poängen
med att koden i lösningsförslagen är den kod som körs: att ett fönster
inte ritas förrän händelseslingan får köra ser du själv genom att fråga
fönstret med \mintinline{python}{winfo_ismapped}
(\cref{sec:clock}), och att nedräkningen gör rätt vid en hel minut och
vid noll ser du i körningen av \texttt{countdown.py}
(\cref{sec:countdown}).  Några påståenden handlar om vad
\mintinline{python}{tkinter} \emph{gör}, och för dem är Pythons egen
dokumentation den rätta källan; adressen står i sammanfattningens fotnot,
med datum.

Men några påståenden kan ingen körning avgöra:
\begin{enumerate}
  \item att det hjälper att försöka själv innan man läser
    lösningsförslaget --- och att det som hjälper är jämförelsen mellan
    det egna försöket och förslaget, inte väntan i sig (övningens
    inledning);
  \item att det är värt något att inte upprepa sig i ett program ---
    principen \emph{DRY}, som materialet åberopar både när den används
    (\cref{sec:timer}) och när den inte används
    (\cref{sec:rect}) --- och att den kommer från Hunt och Thomas;
  \item att det som varierar mot en invariant bakgrund är det som kan
    urskiljas, den variationsteoretiska premiss som lärarnoterna i
    marginalen bygger på;
  \item de missuppfattningar lärarnoterna namnger --- att
    \mintinline{python}{mainloop} skulle betyda \enquote{visa fönstret},
    att en komponent \enquote{är} sitt värde, att parenteserna efter ett
    metodnamn lämnar över metoden --- och att de är \emph{vanliga};
  \item att ett program blir lättare att felsöka om logiken kan köras
    utan gränssnittet (\cref{sec:countdown}).
\end{enumerate}
Hur vet man om sådant stämmer?  Samma fråga möter dig varje gång du läser
en lärobok, en webbsida eller ett svar från en språkmodell --- och varje
gång du själv ska skriva en rapport.  Den här bilagan sammanfattar
arbetsgången.  De två första påståendena är redan prövade i andra
föreläsningar, och \cref{tab:graphics-pastaenden} anger var; de tre
sista är ännu inte prövade, och \cref{sec:graphics-fortsatt} säger vad
som återstår.  Arbetsgången beskrivs utförligare, och genomförd på
flera påståenden, i bilagorna till föreläsningarna \emph{Algoritmiskt
tänkande} och \emph{Upprepningar}.

\begin{table}[htbp]
  \centering
  \caption{Övningens påståenden som frågor, och var svaren finns.  De
    två första är besvarade i andra föreläsningars bilagor; de tre sista
    är öppna, och listade i \cref{sec:graphics-fortsatt}.}
  \label{tab:graphics-pastaenden}
  \begin{tabular}{@{}p{0.48\linewidth}p{0.42\linewidth}@{}}
    \toprule
    Fråga & Svar \\
    \midrule
    Hjälper det att försöka själv före förklaringen?
      & Ja, med måttlig effekt, och effekten kommer av jämförelsen och
        inte av dröjsmålet: belagt i \emph{Upprepningar}, bilaga~C \\
    Är det värt något att inte upprepa sig, och kommer \emph{DRY} från
    Hunt och Thomas?
      & Ja på båda, men \emph{DRY} är en tumregel med gränser: belagt i
        \emph{Algoritmiskt tänkande}, bilaga~E \\
    Är det som varierar mot en invariant bakgrund det som kan urskiljas?
      & Premiss från variationsteorin, med källor i lärarnoternas
        källfil; samma i alla kursens föreläsningar \\
    Är de missuppfattningar lärarnoterna namnger vanliga?
      & Obesvarad: ingen sökning gjord \\
    Blir ett program lättare att felsöka om logiken kan köras utan
    gränssnittet?
      & Obesvarad: ingen sökning gjord \\
    \bottomrule
  \end{tabular}
\end{table}

\section{Gör påståendet till en fråga}

Ett påstående går att kontrollera först när det är tydligt vad som skulle
kunna visa att det är \emph{fel}.  Formulera det därför som en fråga med
ett svar.  \enquote{Kommer \emph{DRY} från Hunt och Thomas?} kan
besvaras med ja eller nej; \enquote{\emph{DRY} är viktigt} kan det inte,
eftersom \enquote{viktigt} är ett omdöme och inte ett faktum.  Ofta
döljer sig flera påståenden i ett: \enquote{en bra princip som kommer
från Hunt och Thomas} är två frågor --- \emph{ursprung} och \emph{värde}
--- och de kräver olika slags källor.

\section{Välj rätt sorts källa}

Olika frågor besvaras av olika källor.
\begin{description}
  \item[Uppslagsverk] för vad ett ord \emph{betyder}.  För vad ett
    programbibliotek \emph{gör} är bibliotekets egen dokumentation
    motsvarande källa: att \mintinline{python}{after} lägger en
    beställning i händelseslingan står i Pythons dokumentation, inte i
    forskningen.
  \item[Primärkällan] för var en idé \emph{kommer ifrån}.
  \item[Översikter] (\foreignlanguage{english}{reviews}) för om något är
    \emph{etablerat}, eller för vad forskningen \emph{sammantaget}
    säger.  Frågan om det hjälper att försöka själv först är just en
    sådan fråga, och den besvarades med en metaanalys.
  \item[Enskilda studier] för ett specifikt resultat --- med förbehållet
    att en enskild studie bara visar att något observerats en gång, i en
    viss population.
\end{description}
Var man hittar dem: universitetsbibliotekets söktjänst, och databaser som
Scopus, Web of Science, IEEE Xplore, DBLP (datalogi) och OpenAlex
(öppen, och trots namnet inte detsamma som öppen tillgång).

\section{Sök så att du kan ha fel}

Den som söker på namn hon redan känner till hittar det hon väntar sig.
Sök därför \emph{ämnesdrivet} --- på begrepp, inte författare.  Några
praktiska regler:
\begin{itemize}
  \item Håll frågorna korta.  Flera korta frågor slår en lång.
  \item Ställ \emph{samma} frågor i alla databaser.
  \item Sök begreppets alla namn.  Det namn som används bland
    programmerare är sällan det som används i litteraturen.
  \item Sök också efter \emph{invändningar}: lägg till ord som
    \enquote{critique}, \enquote{limitations} eller \enquote{harmful}.
    Ett påstående som överlevt en motsökning är värt mer än ett som
    aldrig prövats, och en invändning som hittas blir ett
    \emph{förbehåll} i svaret, inte ett nederlag.
  \item Räkna med att databasernas rankning missar somligt.  Då hämtas
    källan som \emph{känt dokument}: man vet att den finns och slår upp
    den direkt.  Det är i sin ordning, bara det sägs.
\end{itemize}

\section{Läs i fulltext, och skriv ner hur}

En träff i en databas är inte ett belägg.  Sammanfattningen
(\foreignlanguage{english}{abstract}) räcker inte heller: den säger vad
författarna vill ha sagt, inte vad de visat.  Läs hela texten och leta
upp \emph{det ställe} som styrker påståendet --- ett ordagrant citat med
sidnummer.  Räkna aldrig en källa du inte läst.

Skriv sedan ner hur du gick till väga, så att någon annan kan följa
spåret.  I det här materialet står den anteckningen som en kommentar
ovanför varje post i litteraturlistans källfil, med fasta fält:
\begin{description}
  \item[\texttt{CLAIM}] vilket påstående källan ska styrka;
  \item[\texttt{FOUND-VIA}] hur den hittades --- sökfråga och databas,
    eller \enquote{känt dokument};
  \item[\texttt{PICKED}] varför just den, framför andra träffar;
  \item[\texttt{QUOTE}] det ordagranna citatet, med sida;
  \item[\texttt{VERIFIED}] att fulltexten lästs, och var;
  \item[\texttt{COUNTER}] motsökningen och vad den gav;
  \item[\texttt{DATE}] när.
\end{description}
Det tar ett par minuter per källa, och det är det enda som gör det
möjligt att om ett år svara på frågan \enquote{hur vet du det?}

\section{Dra slutsatsen --- med förbehåll}

Svaret skrivs ut tillsammans med det som talar emot.  Att inte upprepa
sig är, enligt bilagan i \emph{Algoritmiskt tänkande}, en tumregel med
undantag och inte en regel utan undantag: upprepningen är risken, men
att undvika den till varje pris är också en risk.  Det är därför
övningen både använder principen, när de tre fälten görs av en och samma
metod (\cref{sec:timer}), och avstår från den, när de två
formknapparna skrivs ut en och en (\cref{sec:rect}).  Och att det
hjälper att försöka själv först gäller, enligt bilagan i
\emph{Upprepningar}, med måttlig effekt och med undantag för yngre
elever och för allmänna färdigheter.  Förbehållen är inte svagheter i
svaret.  De \emph{är} svaret, och de avgör hur påståendet får användas.

\section{En anmärkning om språkmodeller}

Där sökningar gjorts i kursens andra föreläsningar användes en
språkmodell för att \emph{sortera} hundratals träffar --- hör träffen
till ämnet, och åt vilket håll pekar den? --- så att gallringen kunde
granskas.  Modellen sorterade, men varje källa som citeras ska läsas och
väljas av en människa, och även sorteringen ska bekräftas av en
människa.  Använd gärna språkmodeller för att hitta och sortera; låt dem
aldrig vara det sista ledet.  De kan göra fel, men det är du själv som
måste stå till svars för innehållet.

\section{Fortsatt arbete}
\label{sec:graphics-fortsatt}

Den här bilagan redovisar ingen egen sökning.  Två av övningens
påståenden är återanvända från andra föreläsningars bilagor, och tre är
öppna.  De tre är följande, och var och en är markerad i källfilen där
den görs:
\begin{enumerate}
  \item Att nybörjare uppfattar utskrift som \enquote{en rad i taget}
    och därför inte väntar sig att utskriften buffras
    (\cref{sec:term-clock}).
  \item Att nybörjare läser \mintinline{python}{mainloop} som
    \enquote{starta programmet} i stället för \enquote{lämna över
    kontrollen} (\cref{sec:clock}), att de läser av komponenten i
    stället för att lagra värdet (\cref{sec:timer}), att de förväxlar
    att anropa en \foreignlanguage{english}{callback} med att lämna över
    den (\cref{sec:colors}), och att de väntar sig att en ny bindning
    ersätter den förra (\cref{sec:rect}).  Materialet säger att dessa
    fel förekommer, och lärarnoterna säger att de är vanliga; det senare
    är obelagt här.  Frågan är två: förekommer de, och hur ofta?
  \item Att ett program blir lättare att felsöka om logiken kan köras
    utan gränssnittet (\cref{sec:countdown}).  Näraliggande är
    principen om ett ansvar per funktion, belagd i föreläsningen
    \emph{Funktioner}, bilaga~B --- men den gäller ansvarsfördelning,
    inte felsökning, så den räcker inte som svar.
\end{enumerate}
För att besvara dem krävs ämnesdrivna sökningar av samma slag som i de
andra föreläsningarnas bilagor, med motsökningar, och för frågan om hur
vanliga felen är krävs dessutom studier gjorda på en population som
liknar den här kursens.
